\begin{table}[H]
\centering
\caption{Pre-Reform Kaitz Indices by Sector, Lithuania 2018}
\label{tab:kaitz}
\begin{threeparttable}
\begin{tabular}{llccc}
\toprule
ISIC4 & Sector & Kaitz Index & Employment (000s) & Mean Wage (EUR) \\
\midrule
I & Accommodation/Food & 0.540 & 34.4 & 740 \\
A & Agriculture & 0.510 & 98.7 & 785 \\
S & Other Services & 0.420 & 32.5 & 953 \\
P & Education & 0.406 & 135.3 & 986 \\
G & Retail/Wholesale & 0.391 & 234.7 & 1024 \\
F & Construction & 0.384 & 102.6 & 1042 \\
Q & Health & 0.376 & 94.3 & 1064 \\
C & Manufacturing & 0.374 & 219.7 & 1069 \\
H & Transport & 0.350 & 98.9 & 1141 \\
L & Real Estate & 0.337 & 15.2 & 1188 \\
O & Public Admin & 0.331 & 78.6 & 1209 \\
K & Finance & 0.227 & 19.1 & 1762 \\
J & ICT & 0.166 & 32.3 & 2409 \\
\bottomrule
\end{tabular}
\begin{tablenotes}[flushleft]
\small
\item \textit{Notes:} Kaitz index = statutory minimum wage (400 EUR in 2018) divided by sector average monthly earnings. Higher values indicate greater minimum wage binding. Sectors ordered by Kaitz index (most binding first). Employment and wage data from ILO STAT.
\end{tablenotes}
\end{threeparttable}
\end{table}
